\section{Pole observable and electroweak scheme separation}
\label{sec:pole}

The observable used for comparison is the pole-mass ratio
\begin{equation}
  \sPole = 1-\frac{M_{W,\rm pole}^2}{M_{Z,\rm pole}^2}.
  \label{eq:spole}
\end{equation}
This ratio is distinct from a running \(\MSbar\) weak angle and from effective weak mixing angles extracted from neutral-current asymmetries.  The manuscript therefore treats Eq.~\eqref{eq:spole} as the primary target and later relates other schemes to it.

A vector-boson pole is defined by the complex pole of the propagator,
\begin{equation}
  s_V=M_{V,\rm pole}^2-iM_{V,\rm pole}\Gamma_{V,\rm pole}.
  \label{eq:complex-pole}
\end{equation}
When experimental inputs are quoted in a Breit--Wigner convention, the conversion to pole parameters must be performed before using Eq.~\eqref{eq:spole}.  This is one of the numerical audit tasks of the project.

The project seed comparison uses
\begin{equation}
  M_W=M_Z\sqrt{1-\sdV}.
  \label{eq:MWfromMZ}
\end{equation}
For the seed value \(M_Z=91.1880\,\GeV\), this gives \(M_W=80.3748\,\GeV\).  These numbers are placeholders inherited from the working draft and must be audited against the current electroweak source set before final submission.  The calculation script in Appendix~\ref{app:numerics} reproduces the value.

The Standard Model parameter literature distinguishes several vacuum and renormalization conventions.  In particular, pure \(\MSbar\) treatments may define the Higgs vacuum expectation value as the minimum of the loop-corrected effective potential in Landau gauge, eliminating tadpole diagrams at the cost of gauge-choice dependence in intermediate parameters \cite{MartinRobertson2025}.  Such scheme details are relevant for comparing a mass-ratio statement with Lagrangian parameters, while Eq.~\eqref{eq:spole} remains the pole-spectrum target.

\paragraph*{Section status.}  The observable is defined.  The numerical input table requires a current electroweak source audit before final claims.
